% Study sheet: Volume II, Chapter 2, Compute Infrastructure.
% Numbers: \val{...} from compute_infrastructure_studysheet_numbers.tex, written by compute_infrastructure_studysheet_numbers.py.
\documentclass{studysheet}
\input{compute_infrastructure_studysheet_numbers}
\sheetvolume{2}
\sheetchapter{2}{Compute Infrastructure}
\sheetsources{Numbers from the mlsysim registry and the chapter's scenarios; floors are idealized bounds; GB is decimal; assumed inputs are marked.}

\begin{document}
% ===================================================================== FRONT
\frontside
\begin{bigidea}
A model that outgrows one accelerator turns a chip problem into a \hl{building problem}. Memory,
communication, power and reliability limits bind at successive scales, and each level's limit becomes
the next level's input. In C³ terms, the hardware sets the compute and communication bounds within which
the fleet coordinates \hl{sustained useful work}.
\end{bigidea}

\begin{leftcol}
\sheetsec{Key ideas}
\begin{ideas}
\idea{The walls cascade.}{Memory binds at the node, power at the rack, communication across the pod,
  reliability across the fleet; the workload decides which binds.}
\idea{Capacity and bandwidth are separate tests.}{Capacity: does state fit? Bandwidth: how fast does
  low-reuse work, such as decode, run?}
\idea{Plan on sustained, not peak.}{A run sustaining a fraction $\EfficiencyHw$ of peak needs
  $1/\EfficiencyHw$ times a peak plan's GPU-hours.}
\idea{The hierarchy places parallelism.}{Tensor parallelism stays on NVLink; pipeline and data parallelism
  cross nodes. Misplacement shows as exposed collective time.}
\idea{Watts and failures are design inputs.}{The hall fixes each rack's power and cooling. Some GPU
  fails every $\MeanTimeToFailure/\DeviceCount$ hours: buy checkpoints, fast restart.}
\idea{Specialize by the binding constraint.}{Specialization bets the operators last: compute for training,
  bandwidth for decode, capacity for embeddings.}
\end{ideas}

\vspace{12pt}
\fcard{Fleet Law: where the hardware enters}%
  {\DistributedStepTime = \under{\frac{\IronLawComputeTime}{\DeviceCount}}{chips} + \under{\CommunicationTime}{links}
   + \under{\SynchronizationTime}{waiting} - \under{\OverlappedTime}{hiding}}%
  {Stragglers and pipeline bubbles add waiting inside the step.}%
  {Failures, checkpoints and maintenance add time between steps.}%
  {}

\vspace{12pt}
\sheetsec{The accelerator spectrum: a bet on stability}
\begin{sheettable}{@{}l l L@{}}
\thead{Chip} & \thead{Flexibility} & \thead{Pays off when}\\ \midrule
GPU (H100) & high & workloads keep changing\\
TPU (v5p) & moderate & matrix-heavy work runs at scale\\
Custom ASIC & low & one stable workload, huge volume\\
\end{sheettable}
\vspace{3pt}
{\fontsize{8.4}{10.4}\selectfont\color{muted}Generality tax: silicon spent on flexibility, not arithmetic.\par}
\end{leftcol}%
\begin{rightcol}
\sheetsec{Napkin rules}
\begin{sheettable}{@{}l c L@{}}
\thead{Quantity} & \thead{Rule} & \thead{H100 example}\\ \midrule
Ridge point & $\RooflineRidgePoint = \PeakRate/\Bandwidth$ & \val{peaktflops}~TFLOP/s $\div$ \val{bandwidthtb}~TB/s $\approx$ \val{ridge}~FLOP/byte\\
Transfer floor & $\DataVolume/\Bandwidth$ & decode: weight bytes over HBM bandwidth\\
Decode intensity & $\BatchSize$ FLOP/byte & $2\Parameters$ FLOPs per token on $2\Parameters$ FP16 bytes, shared by a batch of $\BatchSize$\\
Training state per GPU & $\dfrac{\val{bytesparam}\,\Parameters}{\TensorparallelDegreeTp\PipelineparallelDegreePp}$ bytes & \val{reftp} $\times$ \val{refpp} split of a 175B model: \val{refstate} GB\\
Compute-only time & $\dfrac{6\Parameters\DatasetSize}{\DeviceCount\PeakRate\EfficiencyHw}$ & then $\div\,\ScalingEfficiency$; add checkpoint, failure, maintenance\\
Failure interval & $\MeanTimeBetweenFailures \approx \MeanTimeToFailure/\DeviceCount$ & \val{mttf} h per GPU\\
\end{sheettable}
\vspace{3pt}
{\fontsize{8.4}{10.4}\selectfont\color{muted}Adam: FP16 weights and gradients, FP32 master and two moments;
leftover HBM holds activations. $\Parameters$ parameters,
$\DatasetSize$ tokens, $\DeviceCount$ GPUs, $\PeakRate$ dense FP16 peak, $\EfficiencyHw$ local efficiency,
$\TensorparallelDegreeTp, \PipelineparallelDegreePp$ tensor and pipeline degrees. One H100: \val{capgib} GiB
$=$ \val{capgb} GB.\par}

\vspace{12pt}
\sheetsec{The bandwidth staircase (links per direction)}
{\centering
\begin{tikzpicture}[sheet]
  \pgfmathsetmacro\u{54}
  \draw[dim] (10,0) -- (232,0);
  \foreach \tk/\tl in {1/1,10/10,100/100,1000/{1,000},10000/{10,000}}{
    \pgfmathsetmacro\x{14+log10(\tk)*\u}
    \draw[dim] (\x,0) -- (\x,-2.5);
    \node[small, anchor=north] at (\x,-4) {\tl};}
  \node[small, anchor=north east] at (232,-14) {GB/s, log scale};
  \pgfmathsetmacro\xi{14+log10(\val{posIB})*\u}
  \pgfmathsetmacro\xn{14+log10(\val{posNVL})*\u}
  \pgfmathsetmacro\xh{14+log10(\val{posHBM})*\u}
  \filldraw[fill=accentsoft, draw=accent] (\xi,0) circle (2.4pt);
  \filldraw[fill=accent, draw=accent] (\xn,0) circle (2.4pt);
  \filldraw[fill=accent, draw=accent] (\xh,0) circle (2.4pt);
  \node[note, anchor=south, align=center] at (\xi,5) {InfiniBand\\port \val{ladderib}};
  \node[note, anchor=south, align=center] at (\xn,5) {NVLink\\\val{laddernvlink}};
  \node[note, anchor=south, align=center] at (\xh,5) {HBM\\\val{ladderhbm}};
  \node[hotbox, anchor=north, align=center] at (\xn,-28) {tensor\\parallel};
  \node[accbox, anchor=north, align=center] at (\xi,-28) {pipeline,\\data parallel};
\end{tikzpicture}\par}

\vspace{12pt}
\sheetsec{The power wall: watts per rack}
\begin{sheettable}{@{}L r L@{}}
\thead{Unit} & \thead{Power} & \thead{Cooling}\\ \midrule
Traditional rack & \val{tradkw} kW & ordinary air\\
DGX H100 node, \val{nodegpus} GPUs & \val{nodekw} kW & \\
DGX H100 rack, \val{nodespr} nodes & \val{rackkw} kW & engineered air, in a suitable facility\\
High-density training rack & \val{liquidkw} kW & liquid\\
\end{sheettable}
\vspace{3pt}
{\fontsize{8.4}{10.4}\selectfont\color{muted}A rack holds only the whole nodes its power feed and cooling carry.\par}
\end{rightcol}

% ===================================================================== BACK
\backside
\begin{answers}
\ans{Finish}{The floor stays \val{wgb} GB $\div$ \val{bandwidthtb} TB/s $=$ \val{wtmem} ms; \val{ybatch} tokens per step
  give \val{ybatch} $\div$ \val{wtmem} ms $\approx$ \val{ytok} tokens/s (about \val{yone} at batch one). Intensity
  \val{yint} FLOP/byte is still below \val{ridge}: bandwidth still binds. A ceiling: it ignores key-value cache reads.}
\ans{1}{$\val{qtotalsci}$ FLOPs $\div$ (\val{qgpus} $\times$ \val{peaktflops} TFLOP/s $\times$ \val{qeta}\%
  $=$ $\val{qratesci}$ FLOP/s) $=$ \val{qhours} h, about \val{qdays} days compute-only; $\div$ \val{effone}
  $\to$ \val{qdaysscaled} days before operational overheads.}
\ans{2}{\val{bytesparam} $\times$ \val{stateparams}B $=$ \val{statetotal} GB; $\div$ \val{reftp} $=$ \val{statetp} GB per GPU,
  over \val{capgb} GB. \val{statepp} stages give \val{stategpu} GB with \val{stateroom} GB left for activations:
  \val{stategpus} GPUs per replica.}
\columnbreak
\ans{3}{\val{buffer} GB takes \val{bufnvlink} ms on NVLink but \val{bufib} ms on one port, \val{nvlinkgap}$\times$
  slower, and tensor-parallel exchanges sit on every layer's critical path. Keep the group in one node; let
  pipeline or data parallelism cross. Evidence: exposed collective time grows in the trace while compute
  per layer stays flat.}
\ans{4}{\val{mttf} h $\div$ \val{fleetgpus} $\approx$ \val{fleetmtbf} h (about \val{fleetmtbfmin} minutes) between
  GPU failures on average: routine, not rare. Buy checkpoint storage and fast restart. (b) \val{fleetgpus} $\div$
  \val{rackgpus} $=$ \val{bigracks} racks $\times$ \val{rackkw} kW $\approx$ \val{bigmw} MW. A \val{tradkw} kW rack feeds
  \val{tradnodes} \val{nodekw} kW node, so \val{tradracks} racks: plan power before buying GPUs.}
\end{answers}

\begin{leftcol}
\sheetsec{Worked example}
\begin{practicebox}\boxtitle{Is batch-one decode limited by math or by memory?}
Llama-3 \val{wparams}B, FP16 weights (\val{wgb} GB, more than one H100 holds), one request, one
hypothetical H100-equivalent path (\val{peaktflops} TFLOP/s, \val{bandwidthtb} TB/s), weights only.\par\vspace{5pt}
{\centering
\begin{tikzpicture}[sheet]
  \pgfmathsetmacro\ux{50}\pgfmathsetmacro\uy{22}
  \draw[dim] (20,0) -- (222,0); \draw[dim] (20,0) -- (20,72);
  \foreach \tk/\tl in {1/1,10/10,100/100,1000/{1,000},10000/{10,000}}{
    \pgfmathsetmacro\x{20+log10(\tk)*\ux}
    \draw[dim] (\x,0) -- (\x,-2.5); \node[small, anchor=north] at (\x,-4) {\tl};}
  \node[small, anchor=north east] at (222,-13) {FLOP per byte};
  \node[small, anchor=south west] at (22,73) {TFLOP/s, log};
  \pgfmathsetmacro\bw{\val{bandwidthtb}}
  \pgfmathsetmacro\pk{\val{peaktflops}}
  \pgfmathsetmacro\xr{20+log10(\pk/\bw)*\ux}
  \pgfmathsetmacro\yr{log10(\pk)*\uy}
  \draw[accent, line width=0.8pt] (20,{log10(\bw)*\uy}) -- (\xr,\yr) -- (222,\yr);
  \draw[accent, dashed, line width=0.5pt] (\xr,0) -- (\xr,\yr);
  \node[note, anchor=south] at (\xr,\yr+2) {ridge \val{ridge}};
  \pgfmathsetmacro\xt{20+log10(\val{posTRAIN})*\ux}
  \filldraw[fill=accent, draw=accent] (20,{log10(\bw)*\uy}) circle (2.4pt);
  \filldraw[fill=white, draw=accent] (\xt,\yr) circle (2.2pt);
  \node[note, anchor=north west] at (24,{log10(\bw)*\uy-3}) {decode, batch one: \val{wintensity}};
  \node[note, anchor=north] at (\xt,\yr-4) {training: \val{rtraini}};
\end{tikzpicture}\par}
\vspace{2pt}
\begin{steps}
\item Math: $2\Parameters$ $=$ \val{wgflop} GFLOP per token; $\div$ \val{peaktflops} TFLOP/s $\approx$ \val{wtcomp} ms.
\item Memory: \val{wgb} GB $\div$ \val{bandwidthtb} TB/s $\approx$ \res{\val{wtmem} ms}.
\item Intensity: \val{wgflop} GFLOP $\div$ \val{wgb} GB $\approx$ \res{\val{wintensity} FLOP/byte}, far below \val{ridge};
  memory is \val{wshare}\% of the summed time.
\end{steps}
\takeaway{Below the ridge, compare bandwidth per dollar, not TFLOP/s; reuse moves the point right.}
\end{practicebox}

\vspace{12pt}
\sheetsec{Your turn: finish it}
\begin{practicebox}\boxtitle{The same path, now \val{ybatch} requests at once}
Same Llama-3 \val{wparams}B weights and H100-equivalent path; only the batch changes, from one to \val{ybatch}.\par\vspace{4pt}
\textbf{Given:} each decode step still reads every weight once, and all \val{ybatch} requests share that read.\par
\textbf{You:} find the step floor, the ceiling in tokens per second, and whether the batch crossed the ridge.
\end{practicebox}
\end{leftcol}%
\begin{rightcol}
\sheetsec{Common mistakes}
\mistake{``Eight times the GPUs trains eight times faster.''}%
  {Scaling efficiency falls as exposed communication grows: \val{effone} at \val{gpusone} GPUs, \val{effeight} at
   \val{gpuseight}. That is \val{hwratio}$\times$ the hardware for \val{speedup}$\times$ the speed, with \val{lost}\% of local
   throughput lost.}{}
\mistake{``A custom ASIC beats a GPU on throughput per watt, so it is always the better buy.''}%
  {Its fixed dataflow bets that today's operators last. A chip takes over a year to design, then serves years
   in a facility built to outlast several generations; a workload shift can strand its efficiency advantage.}{}

\vspace{2pt}
\sheetsec{Check yourself}
\begin{qs}
\item Train the \val{qparams}B model on \val{qtokens}B tokens on the \val{qgpus}-GPU reference cluster at
  \val{qeta}\% local efficiency. What is the compute-only time, and what does a scaling efficiency of \val{effone} make it?
\item \assumed{} a \val{stateparams}B model with mixed-precision Adam and eight-way tensor parallelism on
  H100s. Does its training state fit per GPU? What is the shallowest power-of-two pipeline that fits?
\item A team splits one tensor-parallel group across two nodes. Compare a \val{buffer} GB transfer on
  each link. What should cross instead, and what trace evidence shows the problem?
\item \assumed{} a \val{fleetgpus}-GPU training job. (a) How often does some GPU fail, and what must the
  design buy? (b) On DGX H100 racks, how many racks and megawatts? Could a hall of traditional racks host it?
\end{qs}
\end{rightcol}
\end{document}
